\chapter{Estruturas de Lewis, ressonância e VSEPR}\label{ch:b1:lewis-resonance-vsepr}

A molécula de ozônio, \ce{O3}, é formada por três átomos de oxigênio em
uma cadeia angular. Desenhe-a com as regras do Livro 1 e uma das
ligações sai dupla, a outra simples: uma ligação deveria ser curta, a
outra longa. No entanto, a espectroscopia de micro-ondas encontra as
duas ligações \ce{O-O} do ozônio exatamente iguais, com
\qty{127.8}{pm}, entre a ligação dupla de \ce{O2} (\qty{120.8}{pm}) e a
ligação simples do peróxido de hidrogênio (\qty{147.5}{pm}). Nenhum
desenho isolado está certo; dois desenhos juntos estão. Este capítulo
torna preciso o modelo de Lewis --- \omterm{def:b1:lewis-resonance-vsepr:formal-charge}{cargas formais}, octetos incompletos
ou excedidos, ressonância entre várias estruturas --- e depois o usa para
prever as formas das moléculas e os \omterm{def:b1:lewis-resonance-vsepr:dipole}{momentos dipolares} que delas decorrem.
% ledger: geo:O3.rOO, re:O2, geo:H2O2.rOO

\begin{recall}
O Livro 1 (ano 10) descreveu a \omterm{def:b1:lewis-resonance-vsepr:lewis}{ligação covalente} como um par de elétrons
compartilhado, desenhou \omterm{def:b1:lewis-resonance-vsepr:lewis}{estruturas de Lewis} com \omterm{def:b1:lewis-resonance-vsepr:lewis}{pares ligantes} e não
ligantes e previu quatro formas (linear, angular, trigonal, tetraédrica).
A \omterm{def:b1:periodicity:electronegativity}{eletronegatividade} e suas diferenças foram quantificadas no
\cref{ch:b1:periodicity}; os \omterm{def:b1:quantum-numbers:configuration}{elétrons de valência} de um átomo são lidos
em sua configuração (\cref{def:b1:quantum-numbers:configuration}).
\end{recall}

\section{Estruturas de Lewis e cargas formais}

\begin{definition}[Ligação covalente, estrutura de Lewis]\label{def:b1:lewis-resonance-vsepr:lewis}
Uma \emph{ligação covalente}\index{ligacao covalente@ligação covalente} é um par de elétrons
compartilhado por dois átomos, o \emph{par ligante}\index{par ligante}; dois ou
três pares formam uma ligação dupla ou tripla. Um par de \omterm{def:b1:quantum-numbers:configuration}{elétrons de
valência} que pertence a um único átomo é um \emph{par não
ligante}\index{par nao ligante@par não ligante}. Uma
\emph{estrutura de Lewis}\index{estrutura de Lewis} mostra todos os
\omterm{def:b1:quantum-numbers:configuration}{elétrons de valência} de uma molécula ou íon, como ligações (traços entre
átomos) e pares não ligantes (traços ou pares de pontos sobre um átomo).
Pela \emph{regra do
octeto}\index{regra do octeto}, os átomos do \omterm{def:b1:periodicity:table}{período} 2 (C, N, O, F) ficam
cercados, em uma estrutura estável, por quatro pares, oito elétrons; pela
\emph{regra do dueto}\index{regra do dueto}, o hidrogênio, por um par.
\end{definition}

\begin{definition}[Carga formal]\label{def:b1:lewis-resonance-vsepr:formal-charge}
A \emph{carga formal}\index{carga formal} de um átomo em uma \omterm{def:b1:lewis-resonance-vsepr:lewis}{estrutura
de Lewis} é
\[
  q_F = N_v - N_{\text{nl}} - \tfrac12 N_{\text{lig}},
\]
em que $N_v$ é o número de \omterm{def:b1:quantum-numbers:configuration}{elétrons de valência} do átomo livre,
$N_{\text{nl}}$ o número de elétrons de seus \omterm{def:b1:lewis-resonance-vsepr:lewis}{pares não ligantes} e
$N_{\text{lig}}$ o número de elétrons das ligações que ele forma. Ela é
indicada por $\oplus$ ou $\ominus$ ao lado do átomo quando não é nula.
\end{definition}

\begin{proposition}[As cargas formais somam a carga]\label{prop:b1:lewis-resonance-vsepr:formal-sum}
A soma das \omterm{def:b1:lewis-resonance-vsepr:formal-charge}{cargas formais} dos átomos de uma \omterm{def:b1:lewis-resonance-vsepr:lewis}{estrutura de Lewis} é igual
à carga total da molécula ou do íon.
\end{proposition}

\begin{proof}
Somando $q_F$ sobre os átomos, obtém-se $\sum N_v - (\text{elétrons dos pares
não ligantes} + \text{elétrons das ligações})$, porque os elétrons de cada
ligação são contados pela metade em cada um de seus dois átomos. O
parêntese é o número total de \omterm{def:b1:quantum-numbers:configuration}{elétrons de valência} desenhados na
estrutura, que é igual a $\sum N_v$ menos a carga $z$ da espécie. Logo,
$\sum q_F = z$.
\end{proof}

\begin{method}[Desenhar uma estrutura de Lewis]\label{met:b1:lewis-resonance-vsepr:lewis}
Para desenhar a \omterm{def:b1:lewis-resonance-vsepr:lewis}{estrutura de Lewis} de uma molécula ou íon:
\begin{enumerate}
  \item conte os \omterm{def:b1:quantum-numbers:configuration}{elétrons de valência} $N = \sum N_v - z$ e os pares $N/2$;
  \item desenhe o esqueleto: o átomo menos eletronegativo (exceto H)
        costuma ser o central; hidrogênio e flúor são sempre terminais;
  \item complete os octetos dos átomos terminais com \omterm{def:b1:lewis-resonance-vsepr:lewis}{pares não ligantes}
        e depois coloque os pares restantes no átomo central;
  \item se faltar um octeto ao átomo central, transforme \omterm{def:b1:lewis-resonance-vsepr:lewis}{pares não
        ligantes} de seus vizinhos em ligações múltiplas;
  \item calcule as \omterm{def:b1:lewis-resonance-vsepr:formal-charge}{cargas formais}; entre as estruturas possíveis,
        prefira a que tem menos \omterm{def:b1:lewis-resonance-vsepr:formal-charge}{cargas formais}, com as cargas negativas
        sobre os átomos mais eletronegativos.
\end{enumerate}
\end{method}

\begin{example}[O ácido nítrico]\label{ex:b1:lewis-resonance-vsepr:hno3}
\ce{HNO3}: $N = 1 + 5 + 3 \times 6 = 24$ elétrons, 12 pares. Esqueleto:
nitrogênio central, ligado a três oxigênios, um dos quais leva o
hidrogênio. Completar os octetos com ligações simples deixa o nitrogênio
com seis elétrons; um \omterm{def:b1:lewis-resonance-vsepr:lewis}{par não ligante} de um oxigênio vira uma ligação
\ce{N=O}. \omterm{def:b1:lewis-resonance-vsepr:formal-charge}{Cargas formais}: nitrogênio $5 - 0 - 4 = +1$; o oxigênio com
ligação simples e sem hidrogênio, $6 - 6 - 1 = -1$; os outros, 0. O total
é 0, como exige a \cref{prop:b1:lewis-resonance-vsepr:formal-sum}.
\end{example}

\begin{omfigure}
\setchemfig{atom sep=2.4em}
\begin{tikzpicture}[font=\small]
  \node at (0,0) {\large\chemfig{H-\charge{90=\|,270=\|}{O}-\charge{90:2pt=$\scriptstyle\oplus$}{N}(=[:45]\charge{0=\|,90=\|}{O})-[:-45]\charge{0=\|,240=\|,300=\|,45:3pt=$\scriptstyle\ominus$}{O}}};
  \node[below] at (0,-1.5) {ácido nítrico, \ce{HNO3}};
  \node at (5.2,0) {\large\chemfig{\charge{180=\|}{N}~\charge{0=\|}{N}}};
  \node[below] at (5.2,-1.5) {dinitrogênio, \ce{N2}};
  \node at (9.6,0) {\large\chemfig{\charge{135=\|,225=\|,90:2pt=$\scriptstyle\ominus$}{N}=\charge{90:2pt=$\scriptstyle\oplus$}{N}=\charge{45=\|,315=\|}{O}}};
  \node[below] at (9.6,-1.5) {monóxido de dinitrogênio, \ce{N2O}};
\end{tikzpicture}

{\small \omterm{def:b1:lewis-resonance-vsepr:lewis}{Estruturas de Lewis} com \omterm{def:b1:lewis-resonance-vsepr:lewis}{pares não ligantes} e \omterm{def:b1:lewis-resonance-vsepr:formal-charge}{cargas formais}.}
\end{omfigure}

\section{Além do octeto}

\begin{definition}[Moléculas deficientes em elétrons e hipervalentes]\label{def:b1:lewis-resonance-vsepr:hypervalent}
Uma molécula em que um átomo tem menos de oito \omterm{def:b1:quantum-numbers:configuration}{elétrons de valência} é
\emph{deficiente em elétrons}\index{deficiente em eletrons@deficiente em elétrons}: em \ce{BF3}, o boro
tem seis. Uma molécula em que um átomo do \omterm{def:b1:periodicity:table}{período} 3 ou de \omterm{def:b1:periodicity:table}{período}
posterior está cercado por mais de quatro pares é
\emph{hipervalente}\index{hipervalente}: em
\ce{PCl5} o fósforo forma cinco ligações, em \ce{SF6} o enxofre forma seis.
As moléculas com número ímpar de elétrons, como \ce{NO} e \ce{NO2}, não
podem satisfazer a \omterm{def:b1:lewis-resonance-vsepr:lewis}{regra do octeto} em todos os átomos: um elétron fica
desemparelhado.
\end{definition}

\begin{remark}[Por que o período 3 e não o período 2]\label{rem:b1:lewis-resonance-vsepr:hypervalent}
O nitrogênio forma \ce{NF3}, mas nunca \ce{NF5}, enquanto o fósforo forma
tanto \ce{PF3} quanto \ce{PF5}. Um átomo do \omterm{def:b1:periodicity:table}{período} 2 é pequeno demais
para se ligar a mais de quatro vizinhos; um átomo do \omterm{def:b1:periodicity:table}{período} 3 é grande
o bastante para acomodar cinco ou seis. Encontram-se tanto desenhos de
\ce{SO4^{2-}} ou de \ce{H2SO4} com ligações duplas no enxofre (12
elétrons em torno de S) quanto desenhos com ligações simples e \omterm{def:b1:lewis-resonance-vsepr:formal-charge}{cargas
formais} (um octeto); o segundo reflete melhor a distribuição de cargas,
e ambos dão a geometria correta.
\end{remark}

\begin{definition}[Ácido de Lewis, base de Lewis, ligação dativa]\label{def:b1:lewis-resonance-vsepr:lewis-acid}
Um \emph{ácido de Lewis}\index{acido de Lewis@ácido de Lewis} é uma espécie capaz de aceitar
um par de elétrons em um lugar vago de sua camada de valência (\ce{BF3},
\ce{AlCl3}, \ce{H+}, cátions metálicos); uma \emph{base de Lewis}\index{base de Lewis}
é uma espécie com um \omterm{def:b1:lewis-resonance-vsepr:lewis}{par não ligante} que ela pode compartilhar (\ce{NH3},
\ce{H2O}, \ce{F-}). Quando a base cede o par ao ácido, a ligação formada
é uma \emph{ligação dativa}\index{ligacao dativa@ligação dativa}: uma \omterm{def:b1:lewis-resonance-vsepr:lewis}{ligação covalente} cujos
dois elétrons vieram de um mesmo parceiro.
\end{definition}

\begin{example}[Amônia e trifluoreto de boro]\label{ex:b1:lewis-resonance-vsepr:adduct}
\ce{NH3 + BF3 -> H3NBF3}. No aduto, o boro completa seu octeto; o
nitrogênio, que agora compartilha seu \omterm{def:b1:lewis-resonance-vsepr:lewis}{par não ligante}, porta a \omterm{def:b1:lewis-resonance-vsepr:formal-charge}{carga
formal} $+1$ e o boro, $-1$. Uma vez formada, a ligação \ce{N-B} é uma
\omterm{def:b1:lewis-resonance-vsepr:lewis}{ligação covalente} comum.
\end{example}

\begin{omfigure}
\setchemfig{atom sep=2.1em}
\begin{tikzpicture}[font=\small]
  \node at (0,0) {\large\chemfig{H-\charge{90=\|}{N}(-[:240]H)(-[:300]H)}};
  \node at (1.6,0) {$+$};
  \node at (3.1,0) {\large\chemfig{B(-[:90]F)(-[:210]F)(-[:330]F)}};
  \draw[-{Stealth}, thick] (4.3,0) -- (5.3,0);
  \node at (7.4,0) {\large\chemfig{H-\charge{45:1pt=$\scriptstyle\oplus$}{N}(-[:90]H)(-[:270]H)-\charge{45:1pt=$\scriptstyle\ominus$}{B}(-[:90]F)(-[:270]F)-F}};
\end{tikzpicture}

{\small Uma \omterm{def:b1:lewis-resonance-vsepr:lewis-acid}{ligação dativa}: o \omterm{def:b1:lewis-resonance-vsepr:lewis}{par não ligante} da amônia (uma \omterm{def:b1:lewis-resonance-vsepr:lewis-acid}{base de
Lewis}) ocupa o lugar vago do trifluoreto de boro (um \omterm{def:b1:lewis-resonance-vsepr:lewis-acid}{ácido de Lewis}).
Os \omterm{def:b1:lewis-resonance-vsepr:lewis}{pares não ligantes} do flúor não estão desenhados.}
\end{omfigure}

\section{Ressonância}

\begin{definition}[Estruturas de ressonância, híbrido de ressonância]\label{def:b1:lewis-resonance-vsepr:resonance}
Quando os elétrons de uma molécula podem ser distribuídos em várias
\omterm{def:b1:lewis-resonance-vsepr:lewis}{estruturas de Lewis} que diferem apenas pela posição das ligações $\pi$
e dos \omterm{def:b1:lewis-resonance-vsepr:lewis}{pares não ligantes} --- com os núcleos parados no lugar --- cada
uma delas é uma \emph{estrutura de
ressonância}\index{estrutura de ressonancia@estrutura de ressonância} (ou forma mesomérica), ligada
às outras por uma seta de duas pontas $\leftrightarrow$. A molécula real
não é nenhuma delas: ela é o \emph{híbrido de ressonância}\index{hibrido de ressonancia@híbrido de ressonância},
uma estrutura única cuja distribuição eletrônica é uma média ponderada
das distribuições delas.
\end{definition}

\begin{remark}[Uma seta, não um equilíbrio]\label{rem:b1:lewis-resonance-vsepr:arrow}
A seta $\leftrightarrow$ não descreve uma reação: o ozônio não passa de
uma estrutura para a outra. O híbrido é uma única molécula, descrita por
vários desenhos porque um desenho só não basta. Nunca use
$\rightleftharpoons$ entre \omterm{def:b1:lewis-resonance-vsepr:resonance}{estruturas de ressonância}.
\end{remark}

\begin{method}[Escrever e ponderar estruturas de ressonância]\label{met:b1:lewis-resonance-vsepr:resonance}
Para encontrar as \omterm{def:b1:lewis-resonance-vsepr:resonance}{estruturas de ressonância} de uma espécie:
\begin{enumerate}
  \item parta de uma \omterm{def:b1:lewis-resonance-vsepr:lewis}{estrutura de Lewis}; procure um \omterm{def:b1:lewis-resonance-vsepr:lewis}{par não ligante} ou
        uma ligação $\pi$ vizinha de uma ligação $\pi$, de um lugar vago
        ou de uma carga positiva;
  \item desloque pares com setas curvas (um \omterm{def:b1:lewis-resonance-vsepr:lewis}{par não ligante} vira uma
        ligação $\pi$, uma ligação $\pi$ vira um \omterm{def:b1:lewis-resonance-vsepr:lewis}{par não ligante}), sem
        nunca mover um núcleo nem exceder o octeto de um átomo do \omterm{def:b1:periodicity:table}{período} 2;
  \item recalcule as \omterm{def:b1:lewis-resonance-vsepr:formal-charge}{cargas formais};
  \item pondere as estruturas: as mais importantes obedecem à \omterm{def:b1:lewis-resonance-vsepr:lewis}{regra do
        octeto}, têm menos \omterm{def:b1:lewis-resonance-vsepr:formal-charge}{cargas formais} e põem as cargas negativas sobre
        os átomos mais eletronegativos; estruturas equivalentes têm o
        mesmo peso.
\end{enumerate}
\end{method}

\begin{example}[Ozônio, nitrato, benzeno]\label{ex:b1:lewis-resonance-vsepr:examples}
O ozônio tem duas estruturas equivalentes: cada ligação \ce{O-O} é dupla
em uma e simples na outra, de modo que ambas têm ordem de ligação $1.5$
e o mesmo comprimento, \qty{127.8}{pm}, entre \ce{O=O} e \ce{O-O}. O íon
nitrato \ce{NO3-} tem três estruturas equivalentes; cada ligação \ce{N-O}
tem ordem $4/3$. O benzeno, \ce{C6H6}, tem duas estruturas de Kekulé;
suas seis ligações \ce{C-C} são iguais, \qty{139.7}{pm}, entre a ligação
simples do etano (\qty{153.6}{pm}) e a ligação dupla do eteno
(\qty{133.9}{pm}).
% ledger: geo:O3.rOO, geo:C6H6.rCC, geo:C2H6.rCC, geo:C2H4.rCC
\end{example}

\begin{omfigure}
\vspace*{10pt}
\large
\setchemfig{atom sep=2.4em}
\schemestart
\chemfig{\charge{135=\|,225=\|}{O}=[:-30]\charge{270=\|,90:2pt=$\scriptstyle\oplus$}{O}-[:30]\charge{90=\|,0=\|,270=\|,45:5pt=$\scriptstyle\ominus$}{O}}
\arrow{<->}
\chemfig{\charge{90=\|,180=\|,270=\|,135:5pt=$\scriptstyle\ominus$}{O}-[:-30]\charge{270=\|,90:2pt=$\scriptstyle\oplus$}{O}=[:30]\charge{45=\|,315=\|}{O}}
\schemestop

\medskip
\schemestart
\chemfig{*6(=-=-=-)}
\arrow{<->}
\chemfig{*6(-=-=-=)}
\schemestop
\hspace{3em}
\chemfig{\charge{270:2pt=$\scriptstyle\oplus$}{N}(=[:90]\charge{45=\|,135=\|}{O})(-[:210]\charge{150=\|,240=\|,300=\|,90:2pt=$\scriptstyle\ominus$}{O})-[:330]\charge{30=\|,300=\|,240=\|,90:2pt=$\scriptstyle\ominus$}{O}}
\normalsize
\par\vspace{14pt}
{\small Ressonância no ozônio (um \omterm{def:b1:lewis-resonance-vsepr:lewis}{par não ligante} do oxigênio da direita
vira uma ligação $\pi$, enquanto a ligação $\pi$ da esquerda vira um
\omterm{def:b1:lewis-resonance-vsepr:lewis}{par não ligante}) e no benzeno, e uma das
três estruturas equivalentes do íon nitrato, cuja carga é
compartilhada pelos três oxigênios.}
\end{omfigure}

\begin{definition}[Ligação $\sigma$, ligação $\pi$]\label{def:b1:lewis-resonance-vsepr:sigma-pi}
Uma ligação simples é uma \emph{ligação $\sigma$}\index{ligacao sigma@ligação $\sigma$}:
seu par de elétrons fica ao longo do eixo que une os dois núcleos,
resultado da sobreposição frontal de dois orbitais. A segunda e a
terceira ligações de uma ligação múltipla são \emph{ligações
$\pi$}\index{ligacao pi@ligação $\pi$}: seus pares ficam de cada lado
do eixo, resultado da sobreposição lateral de dois orbitais p
perpendiculares a ele. Uma ligação dupla é uma ligação $\sigma$ e uma
$\pi$; uma ligação tripla, uma $\sigma$ e duas $\pi$.
\end{definition}

\begin{omfigure}
\begin{tikzpicture}[font=\small]
  % sigma: two lobes overlapping head-on
  \fill[omDef!35, opacity=0.8] (0,0) ellipse (0.9 and 0.42);
  \fill[omProp!35, opacity=0.8] (1.1,0) ellipse (0.9 and 0.42);
  \fill (0.1,0) circle (1.6pt); \fill (1.0,0) circle (1.6pt);
  \draw[dashed, black!50] (-1.2,0) -- (2.3,0);
  \node[below] at (0.55,-1.05) {$\sigma$: frontal, sobre o eixo};
  % pi: two p orbitals side by side
  \begin{scope}[xshift=5.2cm]
    \foreach \x/\c in {0/omDef, 1.2/omProp} {
      \fill[\c!35, opacity=0.85] (\x,0.45) ellipse (0.32 and 0.45);
      \fill[\c!15, opacity=0.85] (\x,-0.45) ellipse (0.32 and 0.45);
      \draw[\c!70!black] (\x,0.45) ellipse (0.32 and 0.45) (\x,-0.45) ellipse (0.32 and 0.45);
      \fill (\x,0) circle (1.6pt);
    }
    \draw[dashed, black!50] (-0.7,0) -- (1.9,0);
    \draw[omThm, thick, densely dotted] (0.25,0.75) -- (0.95,0.75) (0.25,-0.75) -- (0.95,-0.75);
    \node[below] at (0.6,-1.05) {$\pi$: lateral, acima e abaixo};
  \end{scope}
\end{tikzpicture}

{\small Os dois tipos de sobreposição. Uma ligação $\sigma$ permite a
rotação em torno de seu eixo; uma ligação $\pi$ não, e é por isso que
uma ligação dupla é rígida
(\cref{ch:b1:stereochemistry-in-depth}).}
\end{omfigure}

\section{O modelo VSEPR}

\begin{definition}[Modelo VSEPR, número estérico]\label{def:b1:lewis-resonance-vsepr:vsepr}
No \emph{modelo VSEPR}\index{modelo VSEPR} (repulsão dos pares de
elétrons da camada de valência), os pares em torno de um átomo central A
se repelem e adotam o arranjo que os mantém o mais afastados possível.
Uma molécula é escrita \ce{AX_nE_m}, em que $n$ é o número de átomos
ligados a A (uma ligação múltipla conta uma vez) e $m$ o número de \omterm{def:b1:lewis-resonance-vsepr:lewis}{pares
não ligantes} de A. O \emph{número estérico}\index{numero esterico@número estérico} $n + m$ fixa o
arranjo dos pares: 2 linear, 3 trigonal plano, 4 tetraédrico, 5
bipirâmide trigonal, 6 octaédrico. A forma da molécula é o arranjo
apenas de seus átomos.
\end{definition}

\begin{proposition}[Geometrias VSEPR]\label{prop:b1:lewis-resonance-vsepr:geometries}
O \omterm{def:b1:lewis-resonance-vsepr:vsepr}{modelo VSEPR} prevê as formas e os ângulos ideais da tabela abaixo;
os ângulos medidos são próximos.
\begin{center}
\small
\begin{tabular}{llllc}
\hline
tipo & arranjo dos pares & forma & exemplo & ângulo medido \\
\hline
\ce{AX2} & linear & linear & \ce{CO2} & $180^\circ$ \\
\ce{AX3} & trigonal plano & trigonal plana & \ce{BF3} & $120^\circ$ \\
\ce{AX2E} & trigonal plano & angular & \ce{SO2} & $119.5^\circ$ \\
\ce{AX4} & tetraédrico & tetraédrica & \ce{CH4} & $109.5^\circ$ \\
\ce{AX3E} & tetraédrico & pirâmide trigonal & \ce{NH3} & $106.7^\circ$ \\
\ce{AX2E2} & tetraédrico & angular & \ce{H2O} & $104.5^\circ$ \\
\ce{AX5} & bipirâmide trigonal & bipirâmide trigonal & \ce{PCl5} & $90^\circ$, $120^\circ$ \\
\ce{AX4E} & bipirâmide trigonal & gangorra & \ce{SF4} & $101.6^\circ$, $173.1^\circ$ \\
\ce{AX3E2} & bipirâmide trigonal & em T & \ce{ClF3} & $87.5^\circ$ \\
\ce{AX2E3} & bipirâmide trigonal & linear & \ce{XeF2} & $180^\circ$ \\
\ce{AX6} & octaédrico & octaédrica & \ce{SF6} & $90^\circ$ \\
\ce{AX5E} & octaédrico & pirâmide de base quadrada & \ce{BrF5} & --- \\
\ce{AX4E2} & octaédrico & quadrado plana & \ce{XeF4} & --- \\
\hline
\end{tabular}
\end{center}
% ledger: geo:CO2.aOCO, geo:BF3.aFBF, geo:SO2.aOSO, geo:CH4.aHCH,
% geo:NH3.aHNH, geo:H2O.aHOH, geo:SF6.aFSF, geo:SF4.aFSF2, geo:SF4.aFSF,
% geo:ClF3.aFClF, geo:XeF2.aFXeF
\end{proposition}
\admitted

\begin{proposition}[Os pares não ligantes fecham os ângulos]\label{prop:b1:lewis-resonance-vsepr:lone-pair-angles}
Um \omterm{def:b1:lewis-resonance-vsepr:lewis}{par não ligante}, retido por um único núcleo, se espalha mais que um
\omterm{def:b1:lewis-resonance-vsepr:lewis}{par ligante} e repele com mais força: as repulsões diminuem na ordem \omterm{def:b1:lewis-resonance-vsepr:lewis}{par
não ligante}/\omterm{def:b1:lewis-resonance-vsepr:lewis}{par não ligante} $>$ \omterm{def:b1:lewis-resonance-vsepr:lewis}{par não ligante}/\omterm{def:b1:lewis-resonance-vsepr:lewis}{par ligante} $>$ \omterm{def:b1:lewis-resonance-vsepr:lewis}{par
ligante}/\omterm{def:b1:lewis-resonance-vsepr:lewis}{par ligante}. Por isso, os ângulos entre ligações diminuem à
medida que se acrescentam \omterm{def:b1:lewis-resonance-vsepr:lewis}{pares não ligantes}
($109.5^\circ$ em \ce{CH4}, $106.7^\circ$ em \ce{NH3}, $104.5^\circ$ em
\ce{H2O}), e em uma bipirâmide trigonal os \omterm{def:b1:lewis-resonance-vsepr:lewis}{pares não ligantes} ocupam as
posições equatoriais, onde têm apenas dois vizinhos a $90^\circ$.
% ledger: geo:CH4.aHCH, geo:NH3.aHNH, geo:H2O.aHOH
\end{proposition}
\admitted

\begin{omfigure}
\setchemfig{atom sep=2.6em}
\renewcommand{\arraystretch}{1.2}
\begin{tabular}{ccccc}
\chemfig{O=C=O} &
\chemfig{B(-[:90]F)(-[:210]F)(-[:330]F)} &
\chemfig{C(-[:90]H)(-[:200]H)(<[:280]H)(<:[:335]H)} &
\chemfig{P(-[:90]Cl)(-[:270]Cl)(-[:180]Cl)(<[:320]Cl)(<:[:30]Cl)} &
\chemfig{S(-[:90]F)(-[:270]F)(-[:180]F)(-[:0]F)(<[:225]F)(<:[:45]F)} \\
\small \ce{AX2}, \ce{CO2} & \small \ce{AX3}, \ce{BF3} & \small \ce{AX4}, \ce{CH4} &
\small \ce{AX5}, \ce{PCl5} & \small \ce{AX6}, \ce{SF6} \\[16pt]
\chemfig{\charge{90=\:}{S}(=[:210]O)(=[:330]O)} &
\chemfig{\charge{90=\:}{N}(-[:200]H)(<[:280]H)(<:[:335]H)} &
\chemfig{\charge{60=\:,120=\:}{O}(-[:215]H)(-[:325]H)} &
\chemfig{\charge{0=\:}{S}(-[:90]F)(-[:270]F)(<[:210]F)(<:[:150]F)} &
\chemfig{Xe(-[:0]F)(-[:90]F)(-[:180]F)(-[:270]F)} \\
\small \ce{AX2E}, \ce{SO2} & \small \ce{AX3E}, \ce{NH3} & \small \ce{AX2E2}, \ce{H2O} &
\small \ce{AX4E}, \ce{SF4} & \small \ce{AX4E2}, \ce{XeF4} \\
\end{tabular}

\medskip
{\small Formas previstas pelo VSEPR. Uma cunha cheia aponta para o
leitor, uma cunha tracejada para trás; as ligações simples ficam no plano
da página. Os \omterm{def:b1:lewis-resonance-vsepr:lewis}{pares não ligantes} do átomo central aparecem como pares de
pontos; \ce{XeF4} está desenhado de frente, com seus dois \omterm{def:b1:lewis-resonance-vsepr:lewis}{pares não
ligantes} apontando para o leitor e para trás.}
\end{omfigure}

\begin{method}[Prever uma forma]\label{met:b1:lewis-resonance-vsepr:vsepr}
Para prever a forma em torno de um átomo central:
\begin{enumerate}
  \item desenhe a \omterm{def:b1:lewis-resonance-vsepr:lewis}{estrutura de Lewis} (\cref{met:b1:lewis-resonance-vsepr:lewis});
  \item conte os átomos ligados ao átomo central ($n$) e seus \omterm{def:b1:lewis-resonance-vsepr:lewis}{pares não
        ligantes} ($m$); escreva \ce{AX_nE_m};
  \item leia o arranjo a partir de $n + m$ e a forma a partir das
        posições dos $n$ átomos;
  \item corrija os ângulos ideais: \omterm{def:b1:lewis-resonance-vsepr:lewis}{pares não ligantes} e ligações
        múltiplas ocupam mais espaço que as ligações simples.
\end{enumerate}
\end{method}

\begin{definition}[Hibridização]\label{def:b1:lewis-resonance-vsepr:hybridisation}
A \emph{hibridização}\index{hibridizacao@hibridização} de um átomo é uma
descrição de suas ligações por orbitais orientados nas direções
previstas pelo VSEPR: um átomo de \omterm{def:b1:lewis-resonance-vsepr:vsepr}{número estérico} 4 é dito sp$^3$
(quatro orbitais equivalentes a $109.5^\circ$), de \omterm{def:b1:lewis-resonance-vsepr:vsepr}{número estérico} 3,
sp$^2$ (três orbitais a $120^\circ$ em um plano, restando um orbital p
perpendicular a ele), de \omterm{def:b1:lewis-resonance-vsepr:vsepr}{número estérico} 2, sp (dois orbitais a
$180^\circ$, restando dois orbitais p). Os orbitais p restantes formam
as ligações $\pi$.
\end{definition}

\begin{example}[Os carbonos da química orgânica]\label{ex:b1:lewis-resonance-vsepr:carbon}
No etano, cada carbono é sp$^3$, tetraédrico; no eteno, sp$^2$, com os
seis átomos em um plano e ângulos próximos de $120^\circ$ (medido
\ce{H-C-H}: $117.6^\circ$); no etino, sp, com os quatro átomos em uma
reta. Os comprimentos de ligação diminuem de 153.6 para 133.9 e para
\qty{120.3}{pm} à medida que se acrescentam ligações $\pi$.
% ledger: geo:C2H4.aHCH, geo:C2H6.rCC, geo:C2H4.rCC, geo:C2H2.rCC
\end{example}

\begin{remark}[Uma descrição, não uma causa]\label{rem:b1:lewis-resonance-vsepr:hybrid-limit}
A \omterm{def:b1:lewis-resonance-vsepr:hybridisation}{hibridização} descreve uma geometria já conhecida; ela não a explica,
e falha quanto às energias dos elétrons, que a teoria dos orbitais
moleculares, no volume do 2.º ano, explica. Ela continua sendo a
linguagem cotidiana da química orgânica: um ``carbono sp$^2$'' é um
carbono plano com um orbital p livre para uma ligação $\pi$.
\end{remark}

\section{Moléculas polares: momentos dipolares}

\begin{definition}[Dipolo de ligação, momento dipolar]\label{def:b1:lewis-resonance-vsepr:dipole}
Em uma ligação entre átomos de \omterm{def:b1:periodicity:electronegativity}{eletronegatividades} diferentes, os
elétrons se deslocam em direção ao átomo mais eletronegativo, que porta
uma carga parcial $-\delta e$, enquanto seu parceiro porta $+\delta e$
($0 < \delta < 1$). O
\emph{dipolo de ligação}\index{dipolo de ligacao@dipolo de ligação} é o vetor $\vec\mu = \delta e\,
\vec d$, de norma $\delta e d$, orientado da carga parcial negativa para
a positiva, sendo $\vec d$ o vetor entre elas. O
\emph{momento dipolar}\index{momento dipolar} de uma molécula é a soma
vetorial de seus dipolos de ligação (e das contribuições de seus \omterm{def:b1:lewis-resonance-vsepr:lewis}{pares
não ligantes}); ele é expresso em coulomb-metros ou em debyes, $\qty{1}{\debye} =
\qty{3.336e-30}{C.m}$. Uma molécula com momento dipolar não nulo é uma
\emph{molécula polar}\index{molecula polar@molécula polar}.
% ledger: const:c (1 D = 1e-21/c C.m)
\end{definition}

\begin{remark}[Convenção de sentido]\label{rem:b1:lewis-resonance-vsepr:convention}
A física desenha o vetor dipolo da carga negativa para a positiva, a
convenção usada aqui. Livros de química mais antigos desenham uma seta
com a cauda cortada apontando para a extremidade negativa; só o sentido
difere.
\end{remark}

\begin{proposition}[Dipolo de uma molécula simétrica]\label{prop:b1:lewis-resonance-vsepr:dipole-sum}
Uma molécula \ce{AX_n} sem \omterm{def:b1:lewis-resonance-vsepr:lewis}{pares não ligantes} em A (\ce{AX2} linear,
\ce{AX3} trigonal, \ce{AX4} tetraédrica, \dots) tem \omterm{def:b1:lewis-resonance-vsepr:dipole}{momento dipolar} nulo,
qualquer que seja a polaridade de suas ligações. Uma molécula \ce{AX2}
angular de ângulo $\theta$ cujas ligações têm dipolo $\mu_b$ tem \omterm{def:b1:lewis-resonance-vsepr:dipole}{momento
dipolar}
\[
  \mu = 2\mu_b \cos\frac{\theta}{2} .
\]
\end{proposition}

\begin{proof}
Nas formas simétricas, os vetores de ligação têm a mesma norma e suas
direções estão dispostas simetricamente em torno de A, de modo que sua
soma é nula (para \ce{AX2} linear, eles são opostos; para \ce{AX3} a
$120^\circ$, três vetores de mesma norma somam zero; para \ce{AX4}, a
soma é invariante pelas rotações do tetraedro, logo é nula). Para a
molécula angular, cada vetor de ligação forma o ângulo $\theta/2$ com a
bissetriz; as componentes perpendiculares à bissetriz se anulam e as
componentes ao longo dela se somam: $2\mu_b\cos(\theta/2)$.
\end{proof}

\begin{omfigure}
\begin{tikzpicture}[font=\small, >={Stealth}]
  % water
  \node[aO] (o) at (0,0) {};
  \node[aH] (h1) at (-0.95,-0.73) {};
  \node[aH] (h2) at (0.95,-0.73) {};
  \begin{scope}[on background layer]
    \draw[bond] (o.center) -- (h1.center); \draw[bond] (o.center) -- (h2.center);
  \end{scope}
  \draw[->, omProp, thick] (-0.25,0.15) -- (-0.95,-0.38);
  \draw[->, omProp, thick] (0.25,0.15) -- (0.95,-0.38);
  \draw[->, omThm, very thick] (0,-0.25) -- (0,-1.45);
  \node[omThm, right] at (0.05,-1.25) {$\vec\mu$};
  \node at (0,0.55) {$\delta^-$};
  \node at (-1.25,-1.05) {$\delta^+$}; \node at (1.25,-1.05) {$\delta^+$};
  \node[below] at (0,-1.7) {\ce{H2O}: $\mu = \qty{1.86}{\debye}$};
  % carbon dioxide
  \begin{scope}[xshift=5cm]
    \node[aO] (a) at (-1.1,0) {}; \node[aC] (c) at (0,0) {}; \node[aO] (b) at (1.1,0) {};
    \begin{scope}[on background layer]
      \foreach \d in {-1.6,1.6} {
        \draw[bond, line width=1.1pt] ([yshift=\d pt]a.center) -- ([yshift=\d pt]c.center);
        \draw[bond, line width=1.1pt] ([yshift=\d pt]c.center) -- ([yshift=\d pt]b.center);}
    \end{scope}
    \draw[->, omProp, thick] (-0.75,0.45) -- (-0.2,0.45);
    \draw[->, omProp, thick] (0.75,0.45) -- (0.2,0.45);
    \node[below] at (0,-1.7) {\ce{CO2}: os dipolos de ligação se anulam, $\mu = 0$};
  \end{scope}
\end{tikzpicture}

{\small \omterm{def:b1:lewis-resonance-vsepr:dipole}{Dipolos de ligação} (em laranja, da carga parcial negativa para a
positiva) e sua soma (em vermelho). Na água, angular, eles se somam; no
dióxido de carbono, linear, eles se anulam.}
% ledger: dip:H2O
\end{omfigure}

\begin{example}[Moléculas polares e apolares]\label{ex:b1:lewis-resonance-vsepr:dipoles}
\ce{CO2}, \ce{BF3}, \ce{CH4} e \ce{CCl4} têm \omterm{def:b1:lewis-resonance-vsepr:dipole}{momento dipolar} nulo.
\ce{H2O} (\qty{1.86}{\debye}), \ce{NH3} (\qty{1.48}{\debye}) e
\ce{SO2} (\qty{1.63}{\debye}) são polares. Ao longo da série \ce{CH3Cl},
\ce{CH2Cl2}, \ce{CHCl3}, o momento diminui, 1.87, 1.62, \qty{1.04}{\debye},
até se anular em \ce{CCl4}: os \omterm{def:b1:lewis-resonance-vsepr:dipole}{dipolos de ligação} \ce{C-Cl} se anulam
cada vez mais. Entre os haletos de hidrogênio, ele acompanha a diferença
de \omterm{def:b1:periodicity:electronegativity}{eletronegatividade}: HF 1.83, HCl 1.09, HBr 0.83, HI \qty{0.45}{\debye}.
% ledger: dip:H2O, dip:NH3, dip:SO2, dip:CH3Cl, dip:CH2Cl2, dip:CHCl3,
% dip:CCl4, dip:HF, dip:HCl, dip:HBr, dip:HI
\end{example}

\section{Exercícios}

\begin{exercise}[$\star$]\label{exo:b1:lewis-resonance-vsepr:1}
Desenhe as \omterm{def:b1:lewis-resonance-vsepr:lewis}{estruturas de Lewis} de \ce{H2O}, \ce{NH3}, \ce{CO2}, \ce{HCN},
\ce{CH2O} (metanal) e \ce{NH4+}, e dê as \omterm{def:b1:lewis-resonance-vsepr:formal-charge}{cargas formais}.
\end{exercise}

\begin{exercise}[$\star$]\label{exo:b1:lewis-resonance-vsepr:2}
Dê o tipo \ce{AX_nE_m}, a forma e o ângulo aproximado de
\ce{H2S}, \ce{PCl3}, \ce{BF3}, \ce{SiH4} e \ce{CO3^{2-}} (átomo central
em primeiro lugar em cada fórmula).
\end{exercise}

\begin{exercise}[$\star$]\label{exo:b1:lewis-resonance-vsepr:3}
Quais destas moléculas são polares: \ce{CCl4}, \ce{CH2Cl2}, \ce{BF3},
\ce{NH3}, \ce{CO2}, \ce{SO2}? Justifique a partir das formas.
\end{exercise}

\begin{exercise}[$\star$]\label{exo:b1:lewis-resonance-vsepr:4}
Identifique o \omterm{def:b1:lewis-resonance-vsepr:lewis-acid}{ácido de Lewis} e a \omterm{def:b1:lewis-resonance-vsepr:lewis-acid}{base de Lewis} em \ce{H+ + H2O -> H3O+} e
em \ce{Cu^{2+} + 4NH3 -> [Cu(NH3)4]^{2+}}. Quais ligações são dativas?
\end{exercise}

\begin{exercise}[$\star\star$]\label{exo:b1:lewis-resonance-vsepr:5}
Escreva duas \omterm{def:b1:lewis-resonance-vsepr:resonance}{estruturas de ressonância} do íon etanoato \ce{CH3COO-}. O que
elas preveem para os comprimentos das duas ligações \ce{C-O}? Compare com o
ácido metanoico \ce{HCOOH}, cujas duas ligações \ce{C-O} medem 120.2 e
\qty{134.3}{pm}.
% ledger: geo:H2CO2.rCO, geo:H2CO2.rCO2
\end{exercise}

\begin{exercise}[$\star\star$]\label{exo:b1:lewis-resonance-vsepr:6}
Desenhe a \omterm{def:b1:lewis-resonance-vsepr:lewis}{estrutura de Lewis} do íon sulfato \ce{SO4^{2-}} com quatro
ligações simples \ce{S-O}, calcule as \omterm{def:b1:lewis-resonance-vsepr:formal-charge}{cargas formais} e preveja a
forma. Quantas \omterm{def:b1:lewis-resonance-vsepr:resonance}{estruturas de ressonância} com duas ligações \ce{S=O} e
sem carga no enxofre podem ser escritas?
\end{exercise}

\begin{exercise}[$\star\star$]\label{exo:b1:lewis-resonance-vsepr:7}
Use o VSEPR para prever as formas de \ce{SF4}, \ce{ClF3}, \ce{XeF2} e
\ce{XeF4}. Explique por que os \omterm{def:b1:lewis-resonance-vsepr:lewis}{pares não ligantes} de uma bipirâmide
trigonal ficam no plano equatorial.
\end{exercise}

\begin{exercise}[$\star\star$]\label{exo:b1:lewis-resonance-vsepr:8}
Os ângulos \ce{H-X-H} são $104.5^\circ$ em \ce{H2O} e $92.1^\circ$ em
\ce{H2S}; $106.7^\circ$ em \ce{NH3} e $93.3^\circ$ em \ce{PH3}. Proponha
uma explicação em termos do tamanho e da \omterm{def:b1:periodicity:electronegativity}{eletronegatividade} do átomo
central.
% ledger: geo:H2O.aHOH, geo:H2S.aHSH, geo:NH3.aHNH, geo:PH3.aHPH
\end{exercise}

\begin{exercise}[$\star\star$]\label{exo:b1:lewis-resonance-vsepr:9}
Dê a \omterm{def:b1:lewis-resonance-vsepr:hybridisation}{hibridização} de cada carbono e do nitrogênio em
\ce{CH3-C#N} (etanonitrila) e em \ce{CH2=CH-CH3}, e o número de
ligações $\sigma$ e $\pi$ de cada molécula.
\end{exercise}

\begin{exercise}[$\star\star\star$]\label{exo:b1:lewis-resonance-vsepr:10}
O \omterm{def:b1:lewis-resonance-vsepr:dipole}{momento dipolar} de \ce{NF3} é de apenas \qty{0.24}{\debye}, contra
\qty{1.48}{\debye} para \ce{NH3}, embora a ligação \ce{N-F} seja mais
polar que a ligação \ce{N-H}. Usando os \omterm{def:b1:lewis-resonance-vsepr:dipole}{dipolos de ligação} e o \omterm{def:b1:lewis-resonance-vsepr:lewis}{par não
ligante} do nitrogênio, explique a diferença.
% ledger: dip:NF3, dip:NH3
\end{exercise}

\begin{exercise}[$\star\star\star$]\label{exo:b1:lewis-resonance-vsepr:11}
A molécula de \ce{SO2} tem ângulo de $119.5^\circ$ e \omterm{def:b1:lewis-resonance-vsepr:dipole}{momento dipolar}
de \qty{1.63}{\debye}. Calcule o dipolo de uma ligação \ce{S-O}; com o
comprimento de ligação de \qty{143.2}{pm}, deduza a carga parcial $\delta$
de cada oxigênio.
% ledger: geo:SO2.aOSO, geo:SO2.rSO, dip:SO2
\end{exercise}

\begin{exercise}[$\star\star\star$]\label{exo:b1:lewis-resonance-vsepr:12}
Desenhe as \omterm{def:b1:lewis-resonance-vsepr:lewis}{estruturas de Lewis} mais importantes do monóxido de
dinitrogênio \ce{N2O} (linear, \ce{N-N-O}) e calcule as \omterm{def:b1:lewis-resonance-vsepr:formal-charge}{cargas formais}.
Os comprimentos de ligação medidos são \ce{N-N} \qty{112.8}{pm} e
\ce{N-O} \qty{118.4}{pm}; em \ce{N2}, a ligação tripla mede
\qty{109.8}{pm}. Qual estrutura pesa mais?
% ledger: geo:N2O.rNN, geo:N2O.rNO, re:N2
\end{exercise}

\section{Problema: As moléculas de um raio}

\begin{problem}[{Problema de fim de semana --- as espécies de nitrogênio e
de oxigênio formadas quando um raio aquece o ar: estruturas de Lewis,
ressonância, formas e as cargas parciais da água}]
\label{pb:b1:lewis-resonance-vsepr:1}
Um raio aquece o ar a milhares de graus: \ce{N2} e \ce{O2}
dão \ce{NO}, depois \ce{NO2}, \ce{N2O}, ozônio e, por fim, ácido nítrico
na chuva. Dados: $\qty{1}{\debye} = \qty{3.336e-30}{C.m}$, $e =
\qty{1.602e-19}{C}$. Os valores medidos são dados quando necessário.

\medskip
\noindent\textbf{Parte I --- \omterm{def:b1:lewis-resonance-vsepr:lewis}{Estruturas de Lewis}.}
\begin{enumerate}
  \item Conte os \omterm{def:b1:quantum-numbers:configuration}{elétrons de valência} de \ce{N2}, \ce{NO}, \ce{NO2},
        \ce{N2O}, \ce{O3} e \ce{HNO3}.
  \item Desenhe a \omterm{def:b1:lewis-resonance-vsepr:lewis}{estrutura de Lewis} de \ce{N2}.
  \item Desenhe uma \omterm{def:b1:lewis-resonance-vsepr:lewis}{estrutura de Lewis} de \ce{NO}. Por que a \omterm{def:b1:lewis-resonance-vsepr:lewis}{regra do
        octeto} não pode ser satisfeita nos dois átomos?
  \item Desenhe duas \omterm{def:b1:lewis-resonance-vsepr:lewis}{estruturas de Lewis} de \ce{N2O}, com \ce{N=N=O} e com
        \ce{N#N-O}, e dê as \omterm{def:b1:lewis-resonance-vsepr:formal-charge}{cargas formais}.
  \item Desenhe uma \omterm{def:b1:lewis-resonance-vsepr:lewis}{estrutura de Lewis} do ozônio e dê as \omterm{def:b1:lewis-resonance-vsepr:formal-charge}{cargas formais}.
  \item Desenhe uma \omterm{def:b1:lewis-resonance-vsepr:lewis}{estrutura de Lewis} de \ce{HNO3} e dê as \omterm{def:b1:lewis-resonance-vsepr:formal-charge}{cargas formais}.
  \item Quais das seis espécies têm um \omterm{def:b1:quantum-numbers:unpaired}{elétron desemparelhado}?
\end{enumerate}

\medskip
\noindent\textbf{Parte II --- Ressonância e comprimentos de ligação.}
\begin{enumerate}[resume]
  \item Escreva as duas \omterm{def:b1:lewis-resonance-vsepr:resonance}{estruturas de ressonância} do ozônio. Que ordem de
        ligação elas preveem?
  \item A ligação \ce{O-O} mede \qty{127.8}{pm} no ozônio,
        \qty{120.8}{pm} em \ce{O2} e \qty{147.5}{pm} em \ce{H2O2}. A
        previsão se confirma?
  \item Escreva as três \omterm{def:b1:lewis-resonance-vsepr:resonance}{estruturas de ressonância} de \ce{NO3-} e dê a
        ordem de cada ligação \ce{N-O}.
  \item Em \ce{HNO3}, as três ligações \ce{N-O} medem 140.6, 121.1 e
        \qty{119.9}{pm}. Atribua-as e explique usando a ressonância.
  \item Explique por que o ângulo \ce{O-N-O} de \ce{NO2} ($134^\circ$) é
        maior que $120^\circ$.
  \item Qual das duas estruturas da questão 4 põe a \omterm{def:b1:lewis-resonance-vsepr:formal-charge}{carga formal} negativa
        no átomo mais eletronegativo?
\end{enumerate}

\medskip
\noindent\textbf{Parte III --- Formas e polaridade.}
\begin{enumerate}[resume]
  \item Dê o tipo \ce{AX_nE_m} do átomo central de \ce{O3},
        \ce{N2O}, \ce{NO3-} e do nitrogênio de \ce{HNO3}, e suas formas.
  \item O ângulo do ozônio é $116.8^\circ$. Explique por que ele é menor
        que $120^\circ$.
  \item O ozônio tem \omterm{def:b1:lewis-resonance-vsepr:dipole}{momento dipolar} de \qty{0.53}{\debye}, o \ce{N2O},
        \qty{0.16}{\debye}, e o \ce{N2}, nenhum. Explique.
  \item Por que \ce{CO2} é apolar enquanto \ce{SO2} é polar?
  \item Preveja o ângulo de \ce{NH3} em relação a $109.5^\circ$ e depois
        compare com o valor medido de $106.7^\circ$.
  \item Mesma pergunta para a água ($104.5^\circ$).
\end{enumerate}
% ledger: geo:O3.rOO, re:O2, geo:H2O2.rOO, geo:HNO3.rNO, geo:HNO3.rNO2,
% geo:HNO3.rNO3, geo:NO2.aONO, geo:O3.aOOO, dip:O3, dip:N2O

\medskip
\noindent\textbf{Parte IV --- As cargas parciais da água.}
O \omterm{def:b1:lewis-resonance-vsepr:dipole}{momento dipolar} da água é \qty{1.857}{\debye}, seu ângulo,
$104.48^\circ$, e o comprimento de sua ligação \ce{O-H}, \qty{95.8}{pm}.
% ledger: dip:H2O, geo:H2O.aHOH, geo:H2O.rOH
\begin{enumerate}[resume]
  \item Expresse o \omterm{def:b1:lewis-resonance-vsepr:dipole}{momento dipolar} da água em função do \omterm{def:b1:lewis-resonance-vsepr:dipole}{dipolo de
        ligação} $\mu_{OH}$ e do ângulo.
  \item Calcule $\mu_{OH}$ em debyes.
  \item Converta-o em coulomb-metros.
  \item Se o \omterm{def:b1:lewis-resonance-vsepr:dipole}{dipolo de ligação} fosse formado por duas cargas inteiras
        $\pm e$ nos dois núcleos, quanto valeria?
  \item Deduza o caráter iônico fracionário $\delta$ da ligação
        \ce{O-H}, a razão entre o \omterm{def:b1:lewis-resonance-vsepr:dipole}{dipolo de ligação} medido e o de cargas
        inteiras.
\end{enumerate}
\end{problem}
